\subsection{Evaluation Metrics and Statistical Analysis}
\label{sec:evaluation_metrics}

Different evaluation metrics were used according to the prediction and diagnostic tasks. For binary anomaly detection, AUPRC was used as the primary metric because the normal and anomalous windows are substantially imbalanced. AUROC was additionally reported to characterize threshold-independent discrimination. For threshold-dependent evaluation, F1, precision, and recall were computed using thresholds selected exclusively on the validation data and then fixed for test evaluation. The same evaluation procedure was applied to MTCL-UAV, without applying point adjustment, to maintain a consistent comparison across detectors.

For fault-domain diagnosis, Macro-F1 was used as the primary metric because it weights the four fault domains equally despite their unequal sample frequencies. Balanced Accuracy, Macro Recall, and per-class precision, recall, and F1 were also reported to characterize class-specific behavior. The same multiclass metrics were used for the complete five-class prediction task, which jointly considers Normal and the four fault domains.

The quality of the diagnostic evidence was evaluated from complementary perspectives. Aggregation conservation error measures whether TreeSHAP contribution mass is preserved when feature-level evidence is successively mapped to telemetry channels, uORB topics, and functional subsystems. Attribution-guided masking measures the decrease in Macro-F1 after masking the most influential features and is compared with random masking of the same size. Consistency@K evaluates whether the top-ranked telemetry evidence agrees with the annotated fault domain, while Hit@K measures whether at least one domain-consistent feature appears among the top-$K$ ranked features.

For software-module inspection, ranking effectiveness was measured using Top-1, Top-3, and Top-5 recall, mean reciprocal rank (MRR), and EXAM score. Lower EXAM values indicate that fewer software modules need to be inspected before reaching the mutation-associated module. Detection recall and ranking latency were additionally recorded in the controlled replay experiments so that conditional ranking quality could be distinguished from the performance of the complete detector-gated pipeline.

Unless otherwise stated, model results were obtained from five random seeds and are reported as mean $\pm$ standard deviation across seeds. All model-selection decisions and operating thresholds were determined from the validation data only. For real-flight evaluations, 95\% confidence intervals were estimated using 1,000 flight-log-level bootstrap repetitions, with complete flight logs rather than overlapping windows used as the resampling unit. This preserves within-flight dependence and avoids treating strongly correlated windows as independent observations.
